% Qwen readout SAE recipe and capacity appendix.

\section{Qwen Readout SAE Recipe and Capacity Sweeps}
\label{app:readout-factorizer-qwen-sweeps}

This appendix gives the Qwen3.5-9B recipe-selection and capacity evidence
summarized in Appendix~\ref{app:readout-factorizer-sweeps}.

\subsection{Qwen3.5-9B Fixed-\texorpdfstring{$k=128$}{k=128} Recipe Selection}
\label{app:readout-factorizer-qwen-recipe}

The first Qwen3.5-9B sweep fixes the interpretability budget at $k=128$ and
selects a no-PCA recipe. A 5k screening sweep identified the recipe family, a
20k continuation evaluated convergence of the selected recipe, and the older 50k
matryoshka run \citep{bussmann2025matryoshka} anchored the comparison against
the previous long-trained baseline. The selected TopK recipe's own convergence ladder provides the
step-count evidence, while the previous recipe serves as an architectural
reference point.

\begin{table*}[!tbp]
\caption{Evidence map for the Qwen3.5-9B fixed-$k=128$ recipe-selection
experiment. Per-run metrics for the 5k screening sweep are provided with
Figure~\ref{fig:app-k-qwen-recipe-selection-frontier}.}
\label{tab:app-k-qwen-recipe-ledger}
\centering
\footnotesize
\setlength{\tabcolsep}{6pt}
\renewcommand{\arraystretch}{1.12}
\begin{tabular}{@{}p{0.20\linewidth}p{0.35\linewidth}p{0.35\linewidth}@{}}
\toprule
Evidence role & Run or run family & What it establishes \\
\midrule
Previous baseline
& Matryoshka TopK, $D=16384$, $k=128$, 50k steps
& Long-trained previous recipe used as an architectural reference: rowEV
$0.511$, top1 $0.807$, KL $0.416$, dead $0.119$, rare $0.236$. \\
\addlinespace
5k selected recipe
& TopK, $D=32768$, $k=128$, row-seeded, hybrid row sampling, delayed prism ramp
& Best recipe-region representative in the screening sweep: rowEV $0.567$,
top1 $0.775$, KL $0.572$, reconstruction error $0.038$, rare $0.060$. \\
\addlinespace
5k screening sweep
& 18 recipe variants varying initialization, architecture, row sampling, prism
schedule, and width
& Supports the selected recipe by matched comparisons: row-seeded
initialization, TopK, delayed ramp, and $D=32768$ are each selected from the
screening grid. \\
\addlinespace
20k selected recipe
& TopK, $D=32768$, $k=128$, row-seeded, hybrid row sampling, delayed prism ramp
& Final reported Qwen3.5-9B fixed-$k=128$ recipe: rowEV $0.621$, top1 $0.846$,
KL $0.296$, reconstruction error $0.027$, dead $0.000$, rare $0.021$. \\
\bottomrule
\end{tabular}
\end{table*}

\begin{table}[!tbp]
\caption{Main effects in the fixed-$k=128$ Qwen3.5-9B recipe sweep.}
\label{tab:app-k-qwen-recipe-effects}
\centering
\footnotesize
\setlength{\tabcolsep}{3pt}
\renewcommand{\arraystretch}{1.08}
\begin{tabularx}{\linewidth}{@{}>{\RaggedRight\arraybackslash}p{0.23\linewidth}Y>{\RaggedRight\arraybackslash}p{0.28\linewidth}@{}}
\toprule
Factor & Evidence & Consequence \\
\midrule
Initialization
& \texttt{row\_seeded} was $2.0\times$ \texttt{random} at 5k
(rowEV $0.563$ vs $0.278$ at $D=32768$).
& Use row-seeded initialization. \\
\addlinespace
Prism-loss schedule
& Delayed ramp kept rowEV unchanged at $0.567$ while reducing selected logit
reconstruction error from $0.064$ to $0.038$ in the matched TopK hybrid cell.
& Ramp $\lambda_{\mathrm{prism}}$ to $10^{-3}$ after $70\%$ of training. \\
\addlinespace
Architecture
& At $D=32768$/5k, plain TopK exceeded matryoshka variants
\citep{bussmann2025matryoshka} in rowEV
($0.567$ vs $0.559/0.554$) at about $3.6\times$ lower cost.
& Use plain TopK for Qwen/Gemma. \\
\addlinespace
Width
& $D=16384$ trailed $D=32768$ by about $0.06$ rowEV in the fixed-$k=128$
recipe sweep.
& Use $D=32768$ for the 9B $k=128$ recipe. \\
\bottomrule
\end{tabularx}
\end{table}

\begin{figure}[!tbp]
    \centering
    \ReadoutSAEGraphic[height=0.30\textheight]{appendix_k_qwen_recipe_selection_frontier.pdf}
    \caption{
        Qwen3.5-9B fixed-$k=128$ recipe-selection frontier. Points: 5k recipe
        variants; black curve: 5k non-dominated frontier. The final 20k TopK
        recipe improves both reconstruction and KL at the fixed
        interpretability budget.
    }
    \label{fig:app-k-qwen-recipe-selection-frontier}
\end{figure}

\ReadoutSAEFloatBarrier

The 9B fixed-$k=128$ recipe is \texttt{topk},
\(D=32768\) (8$\times$), $k=128$, row-seeded
initialization, \texttt{hybrid\_50freq\_50uniform} row sampling, warmup-cosine
learning rate $10^{-3}\rightarrow 10^{-4}$, delayed prism ramp to $10^{-3}$,
20k steps. It reaches rowEV $0.6206$, top1 $0.8457$, KL $0.2955$,
selected logit reconstruction error $0.027$, dead $0.000$, rare $0.021$, and Gini
$0.256$. The 15k$\rightarrow$20k rowEV gain was $+0.0025$, consistent with the
stop rule.

\begin{figure}[!tbp]
    \centering
    \ReadoutSAEGraphic[height=0.28\textheight]{appendix_k_qwen_convergence_ladder.pdf}
    \caption{
        Convergence ladder for the selected Qwen3.5-9B fixed-$k=128$ recipe.
        The final 15k$\rightarrow$20k rowEV gain is small while
        decision-fidelity and usage metrics remain stable, supporting the 20k
        stop rule.
    }
    \label{fig:app-k-qwen-convergence-ladder}
\end{figure}

\paragraph{Recipe-selection conclusion.}
The selected no-PCA TopK recipe improves over the previous long-trained
matryoshka baseline on reconstruction, top-1 agreement, KL, and feature
usage, and its 5k--20k ladder supports the 20k stop rule. The main drivers are
row-seeded initialization, delayed prism ramping, plain TopK, and the
$D=32768$ width choice.

\ReadoutSAEFloatBarrier
\subsection{Qwen3.5-9B Capacity Sweep}
\label{app:readout-factorizer-qwen-capacity}

After recipe selection, we vary capacity along dictionary width and active
sparsity budget. This targeted sweep separates two questions that rowEV alone
can blur: which native $k=128$ readout SAE should be used for fixed-budget
comparisons, and which native $k=256$ readout SAE should serve as the
high-fidelity operating point. Short screening runs selected the larger-budget cells that received
20k training.

A 5k screen selected which larger-capacity cells received 20k training. At
16$\times$, the $k=256$ cell had better rowEV and rare feature
rate than $k=192$ ($0.695/0.100$ vs.\ $0.675/0.164$), despite lower 5k top1,
so it was continued. The 32$\times$, $k=256$ cell had enough rowEV headroom
($0.799$) to justify a width continuation. The native 20k rows below provide
the evidence for the operating-point claims.

\begin{table*}[!tbp]
\caption{Converged 20k native Qwen3.5-9B capacity rows under the selected
recipe family. Native $k=128$ and native $k=256$ rows are separate operating
regimes.}
\label{tab:app-k-qwen-capacity}
\centering
\footnotesize
\setlength{\tabcolsep}{3pt}
\renewcommand{\arraystretch}{1.08}
\begin{tabularx}{\textwidth}{@{}p{0.20\textwidth}ccrrrrY@{}}
\toprule
Run & Width & Eval $k$ & rowEV & top1 & KL & rare & Interpretation \\
\midrule
$D=32768$, train $k=128$
& 8$\times$ & 128 & $0.621$ & $0.846$ & $0.296$ & $0.021$
& strict-budget recipe \\
\addlinespace
$D=65536$, train $k=256$
& 16$\times$ & 256 & $0.761$ & $0.874$ & $0.167$ & $0.061$
& compact $k=256$ comparison \\
\addlinespace
$D=131072$, train $k=256$
& 32$\times$ & 256 & $0.857$ & $0.900$ & $0.105$ & $0.292$
& high rowEV/top1/KL; high rare feature rate \\
\bottomrule
\end{tabularx}
\end{table*}

\begin{figure}[!tbp]
    \centering
\ReadoutSAEGraphic[height=0.28\textheight]{appendix_k_qwen_capacity_usage.pdf}
    \caption{
        Qwen3.5-9B capacity sweep. Increasing width and active sparsity improves
        rowEV and top-1 agreement; rare feature usage is an additional usage
        statistic.
    }
    \label{fig:app-k-qwen-capacity-usage}
\end{figure}

\ReadoutSAEFloatBarrier

Table~\ref{tab:app-k-qwen-capacity} and
Fig.~\ref{fig:app-k-qwen-capacity-usage} show the tradeoff. The
$D=131072$, $k=256$ run has the highest rowEV, top1, and KL metrics in this
Qwen3.5-9B capacity sweep. Its rare feature ladder
$0.300\rightarrow 0.294\rightarrow 0.292$ shows a persistent rare-feature tail,
which is reported as a usage property alongside fidelity metrics.

The 16$\times$, $k=256$ comparison provides the compact high-fidelity point. It
also improves during convergence: from 5k to 20k steps, its rowEV ladder is
$0.719\rightarrow0.749\rightarrow0.758\rightarrow0.761$, and its
rare feature rate falls from $0.100$ to $0.061$.

\ReadoutSAEFloatBarrier
\subsection{Operating-Point Claims}
\label{app:readout-factorizer-operating-point-claims}

The final presentation uses two native operating regimes. The
strict-budget readout uses native $k=128$ readout SAEs for fixed-budget
comparability; after recipe transfer, the selected strict-budget reconstruction
point is Qwen3.5-0.8B 16$\times$
(Appendix~\ref{app:readout-factorizer-small-model-transfer}). The
high-fidelity readout uses native $k=256$ readout SAEs when readout
accuracy is the limiting concern. On Qwen3.5-9B, the 16$\times$, $k=256$ run is
a compact fidelity comparison point with rowEV
$0.761$, top1 $0.874$, KL $0.167$, and rare $0.061$; the 32$\times$, $k=256$
run is the high-capacity Qwen3.5-9B fidelity point with rowEV $0.857$, top1
$0.900$, KL $0.105$, and rare $0.292$.

\begin{figure}[!tbp]
    \centering
    \ReadoutSAEGraphic[width=0.88\linewidth,height=0.26\textheight]{appendix_k_truncation_inversion.pdf}
    \caption{
        Truncating $k=256$-trained dictionaries back to $k=128$ can raise
        rowEV relative to the native $k=128$ recipe while sharply reducing
        top-1 agreement. This rowEV vs.\ fidelity inversion motivates treating
        the fixed interpretability budget as part of the trained object.
    }
    \label{fig:app-k-truncation-inversion}
\end{figure}

\ReadoutSAEFloatBarrier
